\subsection{Reward reweighting: pseudocode for the changes vs the baseline}
\label{app:reward}

Pseudocode for the parts of the GRPO reward that differ between the baseline
(\path{scripts/baseline_scripts/rewards.py}) and the I.3 new-reward run
(\path{scripts/rewards.py}, branch \texttt{reward-reweight}, commit \texttt{e4f6d69}).
The four reward functions are still summed then GRPO-normalised within the group, and
correctness is still gated on the template parsing; only the correctness \emph{shape}
and the term \emph{magnitudes} change, shifting the budget from format ($5.5\to\approx1.0$)
to correctness ($4.5\to5.5$) to counter format-vs-correctness reward hacking.

\begin{verbatim}
# check_answer: coarse ladder (incl. a negative penalty)  ->  smooth graded curve
# BASELINE:
  if   guess == true:                 score = 3.0
  elif guess.strip() == true.strip(): score = 1.5
  else: ratio = guess / true
        score = 0.5 if 0.9<=ratio<=1.1 else 0.25 if 0.8<=ratio<=1.2 else -1.0
        on parse error: score = -0.5
# NEW: graded_correctness(guess, true, MAX=CORRECT_MAX=5.0), shared by both answer terms
  if guess.strip() == true.strip(): return MAX
  err = |guess - true| / max(|true|, 1e-6)        # relative error; parse error -> 0
  return MAX * exp(-CORRECT_K * err)              # CORRECT_K=3.0; smooth, in [0, MAX]

# match_format_exactly:  3.0            ->  FORMAT_EXACT = 0.5   (entry fee)
# match_format_approx.:  +/-0.5 per tag ->  +/-FORMAT_TAG = +/-0.1 per tag (5 tags)
# check_numbers:  1.5 if guess==true else 0  ->  graded_correctness(.., MAX=0.5)
\end{verbatim}
